\documentclass[11pt,a4paper]{article}
\usepackage[utf8]{inputenc}
\usepackage[T1]{fontenc}
\usepackage{amsmath,amssymb,amsthm}
\usepackage{booktabs}
\usepackage{geometry}
\usepackage{hyperref}
\usepackage{listings}
\usepackage{xcolor}

\geometry{margin=2.4cm}

\newcommand{\llbracket}{[\![}
\newcommand{\rrbracket}{]\!]}

\hypersetup{
  colorlinks=true,
  linkcolor=blue!60!black,
  citecolor=blue!60!black,
  urlcolor=blue!60!black
}

\newtheorem{theorem}{Theorem}[section]
\newtheorem{lemma}[theorem]{Lemma}
\newtheorem{definition}[theorem]{Definition}
\newtheorem{invariant}[theorem]{Invariant}
\newtheorem{remark}[theorem]{Remark}

\lstdefinelanguage{Mugda}{
  keywords={sized,data,codata,fun,cofun,let,const,in,eval,mutual,Set,Size},
  sensitive=true,
  morecomment=[l]{--},
  morecomment=[s]{\{-}{-\}},
  basicstyle=\ttfamily\small,
  keywordstyle=\bfseries\color{blue!70!black},
  commentstyle=\itshape\color{gray!70!black},
  columns=fullflexible,
  keepspaces=true,
  frame=single,
  rulecolor=\color{gray!40}
}

\title{\textbf{Mugda Revisited: Sized Types, Size-Change Termination,\\and the Pattern-Refinement Invariant in Lean 4}}
\author{Karl Mehltretter\\
\small Original Diploma Thesis: LMU Munich, December 2007 / April 2008 (Advisor: Andreas Abel)\\
\small Constructive Lean 4 Formalization \& Metatheory, 2026\\
\small \url{https://github.com/kmehltretter82/mugda}}
\date{2026}

\begin{document}
\maketitle

\begin{abstract}
\textbf{Mugda} (2007/2008) was, to our knowledge, the first implemented proof assistant to combine
\emph{Martin-L\"of dependent type theory} and Coquand-style \emph{dependent pattern matching} with
\emph{Lee--Jones--Ben-Amram Size-Change Termination (SCT)} (including the nested matrix semiring $\mathsf{ORDER}_+$)
and \emph{explicit $\mathsf{Size}$ indices} (\texttt{\$ i} for size successor and \texttt{\#} for the closure ordinal $\infty$)
to check both \emph{termination} of inductive definitions (\texttt{fun}) and \emph{productivity} of corecursive definitions (\texttt{cofun}).

Eighteen years later, we revisit Mugda through the lens of transfinite ordinal semantics and machine-checked formalization in \textbf{Lean~4}.
First, we conduct a complete audit of the 2007/2008 system, identifying six soundness and completeness bugs in the interaction between dependent pattern unification, size subtyping, $\mathsf{ORDER}_+$ matrix comparison, and runtime size erasure.
Second, we isolate why naive pattern unification on size indices fails at limit ordinals ($\omega, \omega + 2$): matching $x : \mathsf{D}\;\vec{u}\;\alpha$ against a constructor of target type $\mathsf{D}\;\vec{u}\;(\$ i)$ justifies only the \emph{existential ordinal inequality} $\exists \beta < \alpha,\; \$ \beta \le \alpha$, and \emph{not} the syntactic equality $\alpha = \$ \beta$.
Third, we formulate and prove the \textbf{Pattern-Refinement Invariant} (\texttt{checkNoccEnv} + \texttt{checkNoccDomains} + parameter \texttt{nocc}), showing that substituting $\alpha \mapsto \$ \beta$ during pattern checking is sound across all transfinite ordinals if and only if the refined size variable occurs only in the covariant return type after matching.
Finally, we port the entire Mugda compiler into the \emph{pure constructive subset of Lean~4} (\texttt{Mugda/Core.lean}, with zero \texttt{partial}, zero \texttt{unsafe}, and zero \texttt{sorry}) and verify over 30 semantic and in-place algorithmic metatheorems in \texttt{Mugda/Soundness.lean} and \texttt{Mugda/Metatheory.lean}.
\end{abstract}

\section{Introduction}

Single-argument syntactic guard conditions reject many natural recursive and corecursive programs:
\begin{enumerate}
  \item functions that \emph{permute arguments} (\texttt{add2 x (succ y) = succ (add2 y x)}) or descend lexicographically on nested tuples (\texttt{ack (Pair (succ x) (succ y))}, \texttt{addpair}),
  \item functions that recur on the result of a \emph{size-preserving helper} (Euclidean \texttt{div} and \texttt{gcd} via \texttt{sub}, \texttt{quicksort} via \texttt{pivot}, \texttt{mergesort} via \texttt{split}),
  \item functions that pass a \emph{partially applied recursive call} to a higher-order combinator (\texttt{flat} on rose trees via \texttt{mapL}), and
  \item \emph{corecursive stream and stream-processor definitions} (\texttt{fib}, \texttt{eat}/\texttt{ieat}) where corecursive calls occur under continuations or mutually recursive helpers.
\end{enumerate}

Mugda (Mehltretter, 2007/2008~\cite{mehltretter2007}) unified these paradigms by combining:
\begin{itemize}
  \item \textbf{Syntactic Size-Change Termination in $\mathsf{ORDER}_+$~\cite{lee2001}:} Call extraction constructs call matrices over the recursive order algebra
  \[
    o \in \mathsf{Order} ::= {<} \;\mid\; {\le} \;\mid\; {?} \;\mid\; \mathsf{Mat}(\mathsf{Matrix}\;\mathsf{Order}),
  \]
  computes the idempotent closure under matrix multiplication, and checks that every idempotent call $M = M \cdot M$ has a strictly decreasing diagonal entry (\texttt{isDecr}).
  \item \textbf{Type-Based Termination and Productivity via $\mathsf{Size}$~\cite{abel2006,abel2010}:} Sized inductive families $\mathsf{D}\;\vec{u}\;\alpha$ are covariant in $\alpha$ ($\alpha \le \$ \alpha \le \#$), whereas sized coinductive families $\mathsf{C}\;\vec{u}\;\alpha$ are contravariant in $\alpha$ ($\# \le \$ \alpha \le \alpha$).
\end{itemize}

\section{Transfinite Ordinal Semantics and the Limit-Ordinal Gap}

\subsection{Cumulative Ordinal Interpretation}

In our Lean~4 formalization (\texttt{Mugda/Soundness.lean}), transfinite size ordinals below $\omega \cdot 2$ are represented by the inductive type $\mathsf{SizeOrd}$ with constructors $\mathsf{fin}(n)$ ($n \in \mathbb{N}$) and $\mathsf{omegaPlus}(k)$ ($\omega + k$), extended with the top closure ordinal $\# = \infty$ in $\mathsf{ExtSize}$.
Successor on $\mathsf{ExtSize}$ satisfies $\$ \infty = \infty$ (\texttt{succ\_infty} and \texttt{sinfty\_infty}).

For a monotone operator $F_{\mathsf{D}}$ and a coinductive operator $F_{\mathsf{C}}$, the semantic interpretations at ordinal $\alpha$ are:
\[
  \llbracket \mathsf{D}\;\vec{u}\;\alpha \rrbracket = \bigcup_{\beta < \alpha} F_{\mathsf{D}}\!\left(\vec{u},\; \llbracket \mathsf{D}\;\vec{u}\;\beta \rrbracket\right),
  \qquad
  \llbracket \mathsf{C}\;\vec{u}\;\alpha \rrbracket = \bigcap_{\beta < \alpha} F_{\mathsf{C}}\!\left(\vec{u},\; \llbracket \mathsf{C}\;\vec{u}\;\beta \rrbracket\right).
\]
At the closure ordinal $\infty$, the inductive interpretation stabilizes as a fixed point:
$\llbracket \mathsf{D}\;\vec{u}\;(\$ \infty) \rrbracket = \llbracket \mathsf{D}\;\vec{u}\;\infty \rrbracket$ (proved in \texttt{inOExt\_S\_iff}).
Because Mugda enforces $\infty \not\le \alpha$ for every generic size variable $\alpha$ (\texttt{leqSize VInfty (VGen k) = error}), it avoids the $\infty < \infty$ inconsistency.

\subsection{Why Naive Pattern Unification Fails at Limit Ordinals}

Every constructor $c$ of a sized inductive family $\mathsf{D}$ has a target type of the form $\mathsf{D}\;\vec{u}\;(\$ i)$.
When checking a pattern $c\,i\,\vec{y}$ against an expected domain $\mathsf{D}\;\vec{u}\;\alpha$, Mugda's pattern unifier \texttt{inst} unifies $\alpha$ with $\$ i$, yielding the substitution $\sigma = [\alpha \mapsto \$ i]$.

In 2008 Mugda, $\sigma$ was applied globally across the context $\Gamma$, the remaining argument telescope $\Delta$, and the return type $R(\alpha)$.
However, in the ordinal semantics, $x \in \llbracket \mathsf{D}\;\vec{u}\;\alpha \rrbracket$ guarantees only:
\[
  \exists \beta < \alpha,\quad x \in F_{\mathsf{D}}\!\left(\vec{u},\; \llbracket \mathsf{D}\;\vec{u}\;\beta \rrbracket\right)
  \quad\text{where}\quad \beta < \alpha \implies \$ \beta \le \alpha.
\]
When $\alpha$ is a limit ordinal (e.g.\ $\omega = \mathsf{omegaPlus}(0)$, proved to have no predecessor in \texttt{omega\_is\_limit}), $\alpha \neq \$ \beta$ for all $\beta$.
This gap between the true inequality $\$ \beta \le \alpha$ and the syntactic substitution $\alpha := \$ \beta$ causes two distinct soundness failures:

\begin{enumerate}
  \item \textbf{Deep Sibling Synchronization at $\alpha = \omega + 2$ (Bug 1a, \texttt{cody\_simplified.ma}~\cite{roux2016}):}
  For the Brouwer ordinal type $\mathsf{O} : \mathsf{Size} \to \mathsf{Set}$ with constructors $\mathsf{Z}, \mathsf{S}, \mathsf{L}, \mathsf{M}$, consider at $\alpha = \omega + 2$:
  \[
    t = \mathsf{M}\;(\omega + 1)\;\underbrace{(\mathsf{L}\;\omega\;(\lambda n.\,\mathsf{S}^n\,\mathsf{Z}))}_{\in \llbracket \mathsf{O}\;(\omega + 1) \rrbracket,\; \text{witness } \beta_1 = \omega}\;\underbrace{(\mathsf{S}\;1\;(\mathsf{S}\;0\;\mathsf{Z}))}_{\in \llbracket \mathsf{O}\;2 \rrbracket \subseteq \llbracket \mathsf{O}\;(\omega + 1) \rrbracket,\; \text{witness } \beta_2 = 1}.
  \]
  Although $t \in \llbracket \mathsf{O}\;(\omega + 2) \rrbracket$, matching both children against the deep pattern
  \texttt{M .(\$ (\$ i)) (L .(\$ i) f) (S .(\$ i) (S i z))} forces $\beta_1 = \omega$ to equal $\$ i$ for some $i$, which is false (\textbf{Theorem \texttt{unguarded\_deep\_match\_unsound}} in \texttt{Mugda/Soundness.lean}).

  \item \textbf{Retroactive Context Refinement and Parameter Smuggling (Bug 1b \& Bug 2, \texttt{loop\_wrapped.ma}):}
  If a previous hypothesis in $\Gamma$ has type $f : \mathsf{SNat}\;j \to \bot$ (or is wrapped inside a strictly positive parameter of $\mathsf{Wrap}\;(+ A : \mathsf{Set}) : \mathsf{Size} \to \mathsf{Set}$), matching a subsequent argument $x : \mathsf{SNat}\;j$ against $\mathsf{succ}\;i\;y$ refines $j \mapsto \$ i$ inside $\Gamma$, upgrading $f$ from $\mathsf{SNat}\;j \to \bot$ to $\mathsf{SNat}\;(\$ i) \to \bot$ (\textbf{Theorems \texttt{checkNoccEnv\_tightness} and \texttt{wrap\_not\_covariant\_without\_nocc}}).
\end{enumerate}

\section{The Pattern-Refinement Invariant}

\begin{invariant}[Pattern-Refinement Invariant]
Whenever pattern unification refines a generic size variable $j$ to a successor term $\$ v$ (\texttt{isNewSuccRefinement}) while checking pattern $p$ against domain $A$ in context $\Gamma$ with remaining telescope $\Delta \to R$:
\begin{enumerate}
  \item \textbf{Context Non-Occurrence (\texttt{checkNoccEnv}):} $j$ does not occur free in any type in $\Gamma$.
  \item \textbf{Remaining Domain Non-Occurrence (\texttt{checkNoccDomains}):} $j$ does not occur free in any domain of $\Delta$ (including remaining constructor argument domains).
  \item \textbf{Parameter Non-Occurrence (\texttt{szUsed'} / \texttt{szSizeVarTarget}):} In every sized datatype constructor and admissible function signature, the size variable does not occur in any datatype parameter \texttt{rargs}.
\end{enumerate}
\end{invariant}

\begin{theorem}[Soundness and Tightness of the Pattern-Refinement Invariant]
Under the Pattern-Refinement Invariant, after matching $p$, the refined size variable $j$ occurs only in the return type $R(j)$, which by function admissibility (\texttt{szCheckIndFun}) is covariant in $j$.
For every ordinal $\alpha$ and every witness $\beta < \alpha$, since $\$ \beta \le \alpha$ and $R$ is monotone:
\[
  \llbracket R(\$ \beta) \rrbracket \subseteq \llbracket R(\alpha) \rrbracket.
\]
Hence proving $R(\$ \beta)$ for the constructor's subterms soundly establishes $R(\alpha)$ across all successor and limit ordinals (\textbf{Theorems \texttt{fixed\_pattern\_refinement\_sound}, \texttt{multi\_arg\_pattern\_refinement\_sound}, and \texttt{sized\_recursion\_terminates}}).
Conversely, dropping any of (1), (2), or (3) admits a formal counterexample (\textbf{Theorems \texttt{checkNoccEnv\_tightness}, \texttt{unguarded\_deep\_match\_unsound}, and \texttt{wrap\_not\_covariant\_without\_nocc}}).
\end{theorem}

\section{Algorithmic Verification in Constructive Lean 4}

The entire Mugda compiler is implemented in the pure constructive subset of Lean~4 in \texttt{Mugda/Core.lean} (zero \texttt{partial}, zero \texttt{unsafe}, zero \texttt{sorry}).
Table~\ref{tab:theorems} summarizes the primary machine-checked metatheorems in \texttt{Mugda/Soundness.lean} and \texttt{Mugda/Metatheory.lean}.

\begin{table}[ht]
\centering
\footnotesize
\begin{tabular}{@{}lll@{}}
\toprule
\textbf{Module} & \textbf{Lean 4 Theorem} & \textbf{Verified Property} \\
\midrule
\texttt{Soundness} & \texttt{unguarded\_deep\_match\_unsound} & Limit-ordinal refutation of Bug 1a at $\omega + 2$ \\
\texttt{Soundness} & \texttt{fixed\_pattern\_refinement\_sound} & Single-arg Pattern-Refinement soundness \\
\texttt{Soundness} & \texttt{multi\_arg\_pattern\_refinement\_sound} & Multi-arg telescope refinement soundness \\
\texttt{Soundness} & \texttt{sized\_recursion\_terminates} & Transfinite well-founded sized recursion \\
\texttt{Soundness} & \texttt{inOExt\_S\_iff} & Initial algebra stabilization at $\infty$ (\texttt{\#}) \\
\texttt{Soundness} & \texttt{wrap\_covariant\_of\_nocc} & Covariance of \texttt{Wrap} under parameter \texttt{nocc} \\
\texttt{Soundness} & \texttt{wrap\_not\_covariant\_without\_nocc} & Tightness of Bug 2 parameter \texttt{nocc} \\
\texttt{Soundness} & \texttt{checkNoccEnv\_tightness} & Tightness of \texttt{checkNoccEnv} \\
\texttt{Soundness} & \texttt{inStream\_contravariant}, \texttt{fibFrom\_productive} & Coinductive stream contravariance \& productivity \\
\midrule
\texttt{Metatheory} & \texttt{leqSize\_iterVSucc\_gen/both} & Universal $\$^n$ tower subtyping laws \\
\texttt{Metatheory} & \texttt{inst\_iterVSucc\_infty} & Universal $\$^n$ tower stripping against \texttt{\#} \\
\texttt{Metatheory} & \texttt{matchPat\_eraseSizes\_comm} & Universal size-erasure commutation for \texttt{matchPat} \\
\texttt{Metatheory} & \texttt{succP\_breaks\_size\_erasure} & Tightness of Bug 5 (\texttt{cp (.SuccP \_) = false}) \\
\texttt{Metatheory} & \texttt{nocc\_substVal\_id\_of\_normSizeOrAtom} & Universal \texttt{nocc} $\implies$ \texttt{substVal} identity \\
\texttt{Metatheory} & \texttt{checkNoccEnv\_forall\_mem}, \texttt{checkNoccDomains\_pi} & Universal soundness of \texttt{checkNocc*} \\
\texttt{Metatheory} & \texttt{compScalar\_maxScalar\_distrib\_left/right} & Scalar $\mathsf{Order}$ Kleene-semiring laws \\
\texttt{Metatheory} & \texttt{maxO\_mat\_never\_mat}, \texttt{collapseOrder\_never\_mat} & Bug 3b matrix collapse invariants \\
\texttt{Metatheory} & \texttt{isDecr\_mat\_iff}, \texttt{checkAll\_sound} & Universal SCT diagonal witness characterization \\
\texttt{Metatheory} & \texttt{compScalar\_sound\_sizeOrd}, \texttt{no\_infinite\_decr\_sizeOrd} & Ordinal soundness \& well-foundedness of $\mathsf{Order}$ \\
\texttt{Metatheory} & \texttt{collectCalls\_llet\_includes\_bound} & Bug 4 \texttt{LLet} call-extraction completeness \\
\bottomrule
\end{tabular}
\caption{Summary of machine-checked theorems in \texttt{Mugda/Soundness.lean} and \texttt{Mugda/Metatheory.lean}.}
\label{tab:theorems}
\end{table}

\begin{thebibliography}{9}
\bibitem{mehltretter2007}
Karl Mehltretter.
\newblock \emph{Termination Checking for a Dependently Typed Language}.
\newblock Diploma thesis, LMU Munich, December 2007 (supervised by Andreas Abel).

\bibitem{lee2001}
Chin Soon Lee, Neil D.\ Jones, and Amir M.\ Ben-Amram.
\newblock The Size-Change Principle for Program Termination.
\newblock In \emph{POPL 2001}, pp.\ 81--92. ACM, 2001.

\bibitem{abel2006}
Andreas Abel.
\newblock \emph{A Polymorphic Lambda-Calculus with Sized Higher-Order Types}.
\newblock PhD thesis, LMU Munich, 2006.

\bibitem{abel2010}
Andreas Abel.
\newblock MiniAgda: Integrating Sized and Dependent Types.
\newblock In \emph{PAR 2010}, EPTCS 43, pp.\ 14--28, 2010.

\bibitem{roux2016}
Cody Roux.
\newblock Sized Types and Deep Pattern Matching at Limit Ordinals.
\newblock Agda Issue Tracker (\#1946 and \#2820), 2016/2021.
\end{thebibliography}

\end{document}
